\documentclass[../../../main.tex]{subfiles}

\begin{document}

\section{Cauchy's FE and Its Variants}

Unlike the other problems, if there is one thing that you must
know to solve a FE problem, it would be Cauchy FE. This could
be a very nice problem, but could also be a very mean one.
Before we look at what Cauchy FE is, let's discuss some very
introductory topics from real analysis.

\begin{theorem}
{Density of Rational Numbers}
{Density of Rational Numbers}
For any real numbers $a$ and $b$, it is possible to choose
a rational number $q$ such that $a < q < b$.
\end{theorem}
\begin{proof}
Consider the interval $(na, nb)$ for some integer $n$. Choose a
sufficiently large positive integer $n$ such that $nb - na > 1$.
Because the length of the interval is greater than $1$, we can
choose another integer $m$ such that $na < m < nb$. Therefore,
$a < \frac{m}{n} < b$ and $\frac{m}{n} \in \mathbb{Q}$ since
$m, n \in \mathbb{Z}$.
\end{proof}

This seems quite obvious when you think about it. However,
it is important to consider and question the very seemingly obvious
properties since conditions are very important for Cauchy FE.
Continuing, we could try the same for irrational numbers.

\begin{theorem}
{Density of Irrational Numbers}
{Density of Irrational Numbers}
Given arbitrary real numbers $a$ and $b$, it is possible
to choose a irrational number $q$ such that $a < q < b$.
\end{theorem}
\begin{proof}
First, consider the inequality $\frac{a}{\sqrt{2}} < r <
\frac{b}{\sqrt{2}}$ for some rational number $r$ by Theorem
\ref{theo:Density of Rational Numbers}. Multiplying
$\sqrt{2}$, $a < \sqrt{2} r = q < b$ holds.
\end{proof}

Using these two, we can notice that given two reals, we can
choose a real number between them. Now before moving on to
Cauchy FE, consider the following theorem.

\begin{theorem}
{Bolzano-Weierstrass Theorem}
{Bolzano-Weierstrass Theorem}
Any real number $x$ can be written as the limit of bounded
sequence of rational numbers and of irrational numbers.
\end{theorem}
\begin{proof}
Let $r_n$ be an increasing sequence of rational numbers
defined as $x - \frac{1}{n} < r_n < x - \frac{1}{n+1}$ by
Theorem \ref{theo:Density of Rational Numbers}. Therefore,
\[
x - \frac{1}{n}
< r_n
< x - \frac{1}{n + 1}
< r_{n+1}
\]
holds and $x = \lim_{n \to \infty} r_n$ by Squeeze Theorem.
If $r_n$ is a decreasing sequence defined as
$x + \frac{1}{n+1} < r_n < x + \frac{1}{n}$, then
$r_{n+1} < x + \frac{1}{n+1} < r_n < x + \frac{1}{n}$.
Therefore $x = \lim_{n \to \infty} r_n$ by Squeeze Theorem.
The case of irrational numbers are proven in analogous way.
\end{proof}

Now with these ideas in mind, let's finally discuss
Cauchy's Functional Equations!

\begin{definition}{}{}
Equations with the following form is known as
\textbf{Cauchy's Functional Equations}.
\[
f(x + y) = f(x) + f(y)
\]
\end{definition}

Now this is pretty straight forward for certain input domain.
Let's start by what we can notice just by looking at the
equations. Substituting $x = y = 0$, we get $f(0) = 0$.
Moreover, substituting $y = -x$ gives $f(x) = -f(-x)$ and we
can see that $f$ is an odd function. With these in mind and having
$k = f(1)$, let's start looking into different input domains.

{\Large\textcolor{BrickRed}
{\textbf{I. $\mathbf{f: \mathbb{N} \to \mathbb{R}}$}}}

This case is pretty straightforward. Consider the following
equations for some integer $n$.
\[
f(n)
= f(n - 1) + f(1)
= f(n - 2) + f(1) + f(1)
= \cdots
= n f(1)
\]
Therefore, $f(x) = kx$ if $x$ is defined under natural numbers.
An extension would be using integers.

{\Large\textcolor{BrickRed}
{\textbf{II. $\mathbf{f: \mathbb{Z} \to \mathbb{R}}$}}}

For this case, we will be using the fact that $f$ is an odd function.
For $x > 0$, we have showed that $f(x) = kx$. For $x < 0$,
consider the following equations.
\[
f(x) = -f(-x) = -k(-x) = kx
\]
Therefore with $f(0) = 0$, $f$ is still linear.

{\Large\textcolor{BrickRed}
{\textbf{III. $\mathbf{f: \mathbb{Q} \to \mathbb{R}}$}}}

The next extension that we can have is defining $x$ over the
rational numbers. Let $p$ and $q$ be relatively prime integers.
First consider the following equation.
\[
f \left( \frac{p}{q} \right)
= f \left( \frac{1}{q} + \cdots + \frac{1}{q} \right)
= f \left( \frac{1}{q} \right) \cdot p
\]
Therefore, $f(1) = f \left( \frac{1}{q} \right) \cdot q$ and
substituting back to the equation gives the following.
\[
f \left( \frac{p}{q} \right)
= \frac{f(1)}{q} \cdot p
= \frac{p}{q} k
\]
Therefore, $f(x) = kx$ also holds in this case. The final case
that we can consider is when $f: \mathbb{R} \to \mathbb{R}$.

{\Large\textcolor{BrickRed}
{\textbf{IV. $\mathbf{f: \mathbb{R} \to \mathbb{R}}$}}}

This is the part where it gets quite tricky and where conditions
becomes very important. There are many different conditions that
we could consider for this case, but let's focus on two major
conditions for now.

{\large
\textbf{IV.I $\mathbf{f}$ is an increasing or decreasing function.}}

First, let's consider the case when $f$ is increasing. Assume for
the sake of contradiction that $f(x) \neq kx$. If $f(x) > kx$,
then $x < \frac{f(x)}{k}$ and we can choose a rational $q$
such that $x < q < \frac{f(x)}{k}$ by Theorem \ref{theo:Density of
Rational Numbers}. Therefore, $f(x) \leq f(q) = kq$ since
$q \in \mathbb{Q}$. Continuing, $\frac{f(x)}{k} \leq q <
\frac{f(x)}{k}$ is obtained, which is a contradiction.

This time, assume that $f(x) < kx$. We can choose a rational $q$
such that $\frac{f(x)}{k} < q < x$. Continuing, we get
$kq = f(q) \leq f(x) < kq$ which again is a contradiction.
This proves that if $f$ is an increasing function, then $f(x) = kx$
for $x \in \mathbb{R}$.

The case when $f$ is decreasing is quite simple. Let $g(x) = -f(x)$.
By construction, $g$ is an increasing function and the following
equations hold.
\[
-f(x + y)
= g(x + y)
= g(x) + g(y)
= -f(x) - f(y)
= -(f(x) + f(y))
\]
Therefore, $f(x + y) = f(x) + f(y)$. This completes the proof!

Note that these conditions are necessary. We cannot blindly assume
that $f(x) = kx$ for $x \in \mathbb{R}$ without the condition.
One very, very popular way of knowing if a function is increasing
is $f(xy) = f(x) f(y)$.

If $f(xy) = f(x) f(y)$, then $f(x^2) = f(x)^2 \geq 0$. Therefore
for positive $y$, $f(x) = f(x - y) + f(y) \geq f(x - y)$ and
$f$ is increasing.

\begin{important}
If $f(xy) = f(x) f(y)$, then $f$ is an increasing function.
\end{important}

This is very popular, and will make our lives greatly easier
to use $f(x) = kx$.

{\large\textbf{IV.II $\mathbf{f}$ is continuous.}}

The second very popular condition is when $f$ is continuous.
Note that by Bolzano-Weierstrass Theorem, there exists
an increasing rational sequence $a_n$ and decreasing rational
sequence $b_n$ such that $x = \lim_{n \to \infty} a_n =
\lim_{n \to \infty} b_n$ for any real $x$. Therefore,
the following equations hold.
\begin{align*}
    f(x)
    &= f \left( \lim_{n \to \infty} a_n \right)
    = \lim_{n \to \infty} f(a_n)
    = \lim_{n \to \infty} ka_n
    = kx \\
    &= f \left( \lim_{n \to \infty} b_n \right)
    = \lim_{n \to \infty} f(b_n)
    = \lim_{n \to \infty} kb_n
    = kx
\end{align*}
Therefore, $f(x) = kx$ holds if $f$ is continuous.

This is the basic idea of Cauchy FE! With this in mind,
let's solve few practice problems to see how they are used.





\subsection{Pexider Functional Equation}


\end{document}


